\section{A conformance suite before performance ranking}\label{sec:C22:conformance}
A framework conformance suite should exercise schema validation, state restoration, duplicate dispatch, cancellation, authorization revocation, and trace completeness before comparing task scores. Each fixture has an expected observable contract. For example, a revoked write capability must prevent a new effect even if an old checkpoint contains a permission statement. A committed terminal receipt should be replayed without a new effect; an unresolved dispatch should reconcile before any retry justified by the service's explicit retry contract. A truncated tool result should carry an explicit truncation indication or preserve a recoverable artifact reference.

Conformance is not a universal safety certificate. It establishes behavior for the tested contract and release. A framework can pass deterministic fixtures and still induce poor model decisions through context construction. Therefore the second stage evaluates task outcomes under matched models, tools, and budgets. Separate the effect of the framework from the effect of its recommended prompt, default model, and hosted extensions. When those cannot be separated, the treatment is the complete configured system.

Component effects need not be additive. For two configurable mechanisms, let $\mu_{ij}$ be expected task outcome when mechanism one has state $i\in\{0,1\}$ and mechanism two has state $j\in\{0,1\}$, with the backbone, information access, resource envelope and task population fixed. The interaction is
\begin{equation}
 I=\mu_{11}-\mu_{10}-\mu_{01}+\mu_{00}.
\end{equation}
For example, compaction can change which tool descriptions remain visible, while dynamic discovery changes which descriptions need to survive compaction. A gain from enabling both does not identify either individual contribution; a nonzero $I$ measures their departure from additivity on the declared outcome scale. If an SDK cannot expose all four configurations, the interaction is unidentified from that product comparison. This is an explicit observation limit, not a reason to assign the combined gain to the most prominently advertised feature.

Statistical comparisons should use paired tasks and report uncertainty at the task or session level. Framework adapters can fail only on particular tool schemas, so stratify by schema complexity and effect type. Count startup failures, missing outputs, and dependency outages rather than removing them as setup problems after observing results. A system that requires undocumented manual repair has a materially different operational profile from one that completes unattended under the same fixtures.

Finally, include an implementation-complexity analysis without mistaking it for benchmark quality. Count application-owned components, external dependencies, state stores, custom adapters, and recovery code. These counts do not form a universal simplicity score, but they explain where maintenance responsibility resides. A highly capable managed service may reduce local code while increasing opaque dependencies. A minimal SDK may expose more work but permit a stronger audit. The appropriate choice depends on which uncertainty the deployment can tolerate and which guarantees it must enforce.
